% Generated draft: requires mathematical and source review.
\documentclass[11pt]{article}
\usepackage[margin=25mm]{geometry}
\usepackage{fontspec}
\setmainfont{Latin Modern Roman}
\usepackage{amsmath,amssymb,mathtools,unicode-math}
\setmathfont{Latin Modern Math}
\usepackage{enumitem,needspace,xurl,hyperref}
\setlength{\parindent}{0pt}
\setlength{\parskip}{5pt}
\setlength{\emergencystretch}{3em}
\newcommand{\N}{\mathbb{N}}
\newcommand{\Z}{\mathbb{Z}}
\newcommand{\Q}{\mathbb{Q}}
\newcommand{\R}{\mathbb{R}}
\newcommand{\C}{\mathbb{C}}
\newcommand{\F}{\mathbb{F}}
\newcommand{\A}{\mathbb{A}}
\newcommand{\PP}{\mathbb{P}}
\newcommand{\E}{\mathbb{E}}
\newcommand{\eps}{\varepsilon}
\newcommand{\dd}{\,\mathrm{d}}
\begin{document}
\section*{UnsolvedMath ID: 800087}
\textbf{Source problem code:} catalogue entry
\subsection{Definitions and mathematical statement}
Let \(n\) be a positive integer and let \(G_n=\mathrm{Homeo}_0(S^{2n-1})\), the identity component, in the compact-open topology, of the homeomorphism group of the oriented sphere. Write \(G_n^\delta\) for this group with the discrete topology. The universal oriented flat sphere bundle defines an Euler class
\[
e_{n,\mathbb Z}\in H^{2n}(BG_n^\delta;\mathbb Z),
\qquad e_n\in H^{2n}(BG_n^\delta;\mathbb R),
\]
where \(e_n\) is the image of \(e_{n,\mathbb Z}\) under change of coefficients. Here the discrete-group interpretation is suggested by the source's discussion of flat bundles and holonomy; the source does not explicitly specify the topology or coefficients in its first cohomology expression.

For a discrete group \(\Gamma\), let \(H_b^k(\Gamma;\mathbb R)\) denote the cohomology of bounded real group cochains, with trivial coefficients, and let
\[
c_\Gamma:H_b^k(\Gamma;\mathbb R)\longrightarrow H^k(B\Gamma;\mathbb R)
\]
be the comparison map. An ordinary real cohomology class is called bounded if it lies in the image of this map. A discrete group is amenable if it admits a left-invariant mean on its bounded real-valued functions.

\textbf{Supplied phenomenon.} For \(n>1\), the record reports that the Euler class of flat \(S^{2n-1}\)-bundles is unbounded, while the Hirsch--Thurston result cited there gives vanishing for amenable holonomy. In the real, discrete-group interpretation, these assertions take the form
\[
e_n\notin\operatorname{im}(c_{G_n^\delta}),
\qquad
\operatorname{res}^{G_n^\delta}_{A}(e_n)=0
\quad\text{for every amenable subgroup }A\leq G_n^\delta.
\]
These are supplied background claims, not independently verified assertions of this draft. The coefficient scope of the cited vanishing result is not specified in the record; integral vanishing is therefore not asserted here.

\textbf{Questions from the source.} Explain the coexistence of the reported unboundedness and amenable-subgroup vanishing. In particular, is the Euler class hyperbolic in Gromov's sense, as referenced by the source through \texttt{brunnbauer2024atoroidal}? The record supplies neither a definition of hyperbolicity nor the precise space and realization conditions to which that definition is to be applied. Consequently this part remains a named question with an unresolved formulation qualification.

The source also asks about a nontrivial bounded Euler class for volume-preserving sphere homeomorphisms. To express its possible ordinary-cohomology meaning, fix a volume measure \(\mu\) on \(S^{2n-1}\). The record does not specify that measure or whether its notation \(\mathrm{Homeo}^{\mathrm{vol}}_0(S^{2n-1})\) means
\[
K_\mu=\{g\in G_n:g_*\mu=\mu\}
\quad\text{or}\quad
K_\mu^0=\bigl(\mathrm{Homeo}(S^{2n-1},\mu)\bigr)_0.
\]
Once the intended subgroup \(K\) is specified, put \(e_{n,K}=\operatorname{res}^{G_n^\delta}_{K^\delta}(e_n)\). One precise reading of this question is
\[
e_{n,K}\neq 0
\quad\text{and}\quad
e_{n,K}\in\operatorname{im}\!\left(
 c_{K^\delta}:H_b^{2n}(K^\delta;\mathbb R)
 \longrightarrow H^{2n}(BK^\delta;\mathbb R)
\right).
\]
Another possible reading asks for a nonzero bounded-cohomology class mapping to \(e_{n,K}\), without requiring its ordinary-cohomology image to be nonzero. These readings are distinct, and the supplied wording does not resolve them. The condition \(n>1\) is explicit for the reported unboundedness; its intended scope over the subsequent questions also requires confirmation.
\subsection{Short English statement}
For odd-dimensional spheres, the supplied record asks why the Euler class can be unbounded yet vanish on every amenable subgroup. It asks whether Gromov hyperbolicity explains this behavior and whether restricting to volume-preserving homeomorphisms yields a nontrivial bounded Euler class. The precise hyperbolicity convention, subgroup, coefficient scope, and meaning of nontriviality remain unresolved. This is a statement-only draft requiring human mathematical review.
\subsection{Sources}
\begin{itemize}
\item Local UnsolvedMath registry, supplied complete record with permanent ID \(800087\), exact source problem code \texttt{catalogue entry}, titled ``Foliations and Diffeomorphism Groups --- Question 28.'' The record attributes the problem to Nariman and Yazdi, proposed in 2025, and labels it open. That status is supplied registry metadata, not an independently verified literature assessment.
\item Nariman and Yazdi, \emph{Problem List on Foliations and Diffeomorphism Groups} (2025), Question 28, \url{https://arxiv.org/abs/2503.18274}. The record reports access on 2026-10-16 and extraction by \texttt{arxiv\_latex}; these are supplied provenance details, not a claim of source access by this draft.
\item Hirsch and Thurston, cited in the supplied statement by the reference key \texttt{hirsch1975foliated}, for the reported Euler-class vanishing under amenable holonomy. No complete bibliographic entry, theorem statement, or separate URL is supplied.
\item Reference key \texttt{brunnbauer2024atoroidal}, cited in the supplied statement concerning hyperbolic classes in Gromov's sense. No complete bibliographic entry, definition, or separate URL is supplied.
\end{itemize}
\end{document}
